\chapter[AI Inputs: English and Visual-State Learning]
{Developing the AI Inputs: English Understanding and Visual-State Learning}
\label{chap:ai-development}

\chapterguide{Learned models could assist with task and scene interpretation,
but neither was allowed to issue commands. I therefore built a constrained
English interface and a rail-isolated visual estimator that supply the planner
with two checked records.}

\contributionboundary
  {The starting repository contained no Room~315 language pipeline,
  paired-camera corpus or trained visual-state estimator.}
  {The supervisors supplied platform knowledge and reviewed the methods and
  claims.}
  {I integrated Qwen2.5-1.5B-Instruct \citep{qwen25,qwen25gguf} through
  \code{llama-cpp-python} v0.3.16 and \code{llama.cpp}
  \citep{llamacpppython,llamacpp}. I also reused TorchVision's
  ImageNet-pretrained ResNet-18
  \citep{torchvisionresnet18,resnet,imagenet}; I did not develop these external
  models or tools.}
  {I developed the constrained language contract and dialogue pipeline,
  generated and partitioned the paired-camera corpus, trained the rail-isolated
  V4 estimator from my V3 checkpoint, and integrated its hash-bound runtime.}
  {I ran the pinned ten-case language evaluation, the Validation and
  post-selection regression checks, and the preregistered one-shot Final Test
  reported in \cref{chap:results}.}

\section{Building the English task interface}

\subsection{Local model selection and configuration}

Qwen2.5-1.5B-Instruct is the pretrained model selected for the interface
\citep{qwen25}. In my pipeline, it proposes structured task fields from English
requests. I ran the quantised GGUF file distributed in the official model card
\citep{qwen25gguf} locally on the CPU through \code{llama-cpp-python} v0.3.16,
a Python binding for \code{llama.cpp} \citep{llamacpppython,llamacpp}. I did not
fine-tune the model. The exact deterministic inference configuration and
artefact identity are indexed through \cref{app:reproducibility}.

Without fine-tuning, a versioned prompt constrains the model to one
\code{SemanticParseEnvelope} JavaScript Object Notation (JSON) object over permitted task fields,
identities, rails, stations and slots; it excludes PDDL, plan steps and device
or shuttle commands.

\subsection{From model output to a confirmed task}

A rule-based extractor records explicit facts (for example, shuttle colour,
rail or slot), Qwen proposes intent fields, and fixed-precedence fusion prevents
the proposal from overwriting explicit or confirmed context. The resulting
non-executable \code{TaskGoalDraft} is clarified, domain-validated and confirmed
before an immutable \code{TaskGoal} is issued.

For the indirect request ``Could you get a loaded shuttle over to the right
slot~4?'', the saved record shows the model supplying
\code{goal\_type=transport} and deterministic fusion/validation constraining
side, payload and destination. The ten-case evaluation is reported in
\cref{sec:language-contract-evidence}.

\section{Generating the Room~315 visual corpus}

\subsection{Paired-camera corpus and partitioning}

The generator enumerated all $2^8-1=255$ non-empty subsets of identities
L1--L4/R1--R4 and produced eight geometrically varied scenes per subset by
changing valid rail positions, payload states and switches. After each Gazebo
scene settled, it captured one synchronised image per overhead camera: 2,040
paired observations, or 4,080 RGB files. \Cref{fig:paired-camera-training}
shows one input.

One scene is one labelled paired-camera record; a visible-shuttle target is one
present identity with complete supervision and is the denominator used for the
per-shuttle metrics.

\begin{figure}[htbp]
  \centering
  \begin{subfigure}[t]{0.485\textwidth}
    \includegraphics[width=\textwidth]{f13_left_raw.png}
    \caption{Left-camera RGB input}
  \end{subfigure}\hfill
  \begin{subfigure}[t]{0.485\textwidth}
    \includegraphics[width=\textwidth]{f13_right_raw.png}
    \caption{Right-camera RGB input}
  \end{subfigure}
  \caption[Paired Room~315 camera observation used for visual training.]
  {The two synchronised RGB images form the complete learned input for one
  scenario. Simulator state is used to construct labels in a separate file and
  is never appended to the image tensor.}
  \label{fig:paired-camera-training}
\end{figure}

Gazebo identity, rail coordinate, payload state and projected box supply labels
stored separately from image references. Only paired RGB images are model
features; depth, controller state, point clouds and other oracle metadata remain
label-generation data.

\paragraph{Historical grouped protocol.}
The predecessor V3 experiment kept all variants of a presence configuration in
one partition. Its Validation selected epoch~14, and its Test was consumed once
after selection on 30 July 2026; it is therefore historical evidence, not the
current Final Test. Neither historical 256-scene partition directly enters V4,
although V3's Validation-selected backbone provides indirect initialisation
lineage. Current loaders denylist both historical partitions. Their split counts
are summarised in \cref{app:technical}, while the audit record is indexed through
\cref{app:reproducibility}.

\subsection{Current training, Validation and regression-set roles}

The current pool combines 1,528 legacy grouped-training scenes with 4,000
training-only hard cases. Weighted sampling emphasises rare rail--segment cells,
boundaries, switch zones and multi-blocker scenes. These cases were generated
before Final Test access rather than mined from its errors.

The separate hard-case Validation (512 scenes; 2,254 present targets) selected
V4, drove patience and supplied temperature/segment-threshold fitting. Only
after selection, the 256-scene post-selection regression set (1,085 targets)
checked the frozen candidate for degradation; it affected neither selection nor
calibration and is not Test evidence. Accessing the current Final Test remained prohibited
throughout development; it supplies the principal evidence below.

The labelled partitions were released only after the immutable attempt had
completed; this later publication changes availability, not their experimental
roles. \Cref{app:reproducibility} identifies the release.

\section{Designing and training the visual-state estimator}

The \emph{visual model} is the learned network; the \emph{visual-inference
runtime} adds preprocessing and observation validation; the \emph{closed-loop
runtime} further adds the task gateway, planner, supervisor and controllers.

\paragraph{Evolution of the visual pipeline.}
The early versions primarily evolved the data contract: V1 established
paired-camera inputs while keeping simulator-derived labels outside the model
input, and V2 added continuous along-rail localisation. V3 fixed the
representation to the eight known shuttle identities, resized both cameras to
$224\times224$, concatenated global features from a shared ImageNet-pretrained
ResNet-18 and predicted 200 values through one head. Review of the hard-case
behaviour exposed three mismatches between that contract and the known camera
geometry and topology. \Cref{tab:v3-v4-design-response} maps each diagnosis to
the corresponding V4 decision. V4 is the active visual model used by the
evaluated runtime; its results are in \cref{sec:visual-v4-results}.

\begin{table}[H]
\centering
\small
\renewcommand{\arraystretch}{1.08}
\caption{V3 limitations and the corresponding V4 design responses.}
\label{tab:v3-v4-design-response}
\begin{tabularx}{\textwidth}{@{}p{0.20\textwidth}p{0.31\textwidth}X@{}}
\toprule
\textbf{Design concern} & \textbf{V3 observation} & \textbf{V4 response}\\
\midrule
Image geometry & Squaring the native $4{:}3$ images to $224\times224$ changed
their proportions & Aspect-preserving $320\times240$ preprocessing retains the
camera geometry\\ \tablerowrule
Rail interaction & Concatenated global features allowed one camera to influence
the other rail's prediction & Low-level weights remain shared, but each rail has
an independent late branch and identity queries\\ \tablerowrule
Learned-state scope & One 200-value head also predicted side, metric position
and segment length, although topology already fixes them & The model learns 168
uncertain visual values; identity and topology derive the remaining quantities
deterministically\\
\bottomrule
\end{tabularx}
\end{table}

\subsection{Rail-isolated 168-value learned contract}

Each $4{:}3$ image is resized to $320\times240$, ImageNet-normalised and
processed separately by shared ResNet-18 weights through \code{layer3}. Each
rail then has its own \code{layer4} and four identity-query heads; no camera
feature enters the other rail's late branch.

For each of L1--L4/R1--R4, the model outputs 14 segment scores, two loaded-state
scores, a four-value normalised box and one within-segment ratio: 21 values per
identity and 168 in total. Identity and the decoded topology then determine
side, segment length and metric position; these known quantities are not learned.
The output composition is summarised in \cref{app:technical}; its exact ordering
and diagnostic compatibility adapter remain in the model records indexed by
\cref{app:reproducibility}.

\begin{figure}[p]
  \centering
  \begin{tikzpicture}[
  font=\sffamily\footnotesize,
  execute at begin node={\hyphenpenalty=10000\exhyphenpenalty=10000},
  input/.style={draw=black!42,rounded corners=2pt,fill=black!3,
    align=center,text width=66mm,minimum height=13mm,inner sep=1.4mm},
  encoder/.style={draw=mfjablue,very thick,rounded corners=2pt,
    fill=mfjalight,align=center,text width=66mm,minimum height=24mm,
    inner sep=1.5mm},
  learned/.style={draw=mfjablue,very thick,rounded corners=2pt,
    fill=mfjablue!5,align=center,text width=152mm,minimum height=14mm,
    inner sep=1.5mm},
  contract/.style={draw=black!55,very thick,rounded corners=2pt,
    fill=white,align=center,text width=152mm,minimum height=30mm,
    inner sep=1.6mm},
  runtimeinput/.style={draw=black!48,rounded corners=2pt,fill=black!3,
    align=center,text width=68mm,minimum height=13mm,inner sep=1.3mm},
  gate/.style={draw=mfjagreen,very thick,rounded corners=2pt,
    fill=mfjagreen!7,align=center,text width=152mm,minimum height=19mm,
    inner sep=1.6mm},
  presentcase/.style={draw=mfjagreen,very thick,rounded corners=2pt,
    fill=mfjagreen!7,align=center,text width=47mm,minimum height=23mm,
    inner sep=1.4mm},
  absentcase/.style={draw=black!58,very thick,rounded corners=2pt,
    fill=black!3,align=center,text width=47mm,minimum height=23mm,
    inner sep=1.4mm},
  rejectcase/.style={draw=mfjared,very thick,rounded corners=2pt,
    fill=mfjared!6,align=center,text width=47mm,minimum height=23mm,
    inner sep=1.4mm},
  accepted/.style={draw=mfjagreen,very thick,rounded corners=2pt,
    fill=mfjagreen!10,align=center,text width=152mm,minimum height=23mm,
    inner sep=1.6mm},
  paneltitle/.style={anchor=west,font=\sffamily\small\bfseries,
    text=mfjablue},
  flow/.style={-{Latex[length=2.1mm]},very thick,draw=mfjablue},
  gateflow/.style={-{Latex[length=2.1mm]},very thick,
    draw=mfjagreen!80!black},
  caseflow/.style={-{Latex[length=1.9mm]},very thick,draw=black!75}
]
  % Panel A: the actual paired-view architecture used in training and runtime.
  \node[paneltitle] at (-79mm,11mm)
    {A \enspace LEARNED PAIRED-VIEW STATE ESTIMATOR};
  \draw[mfjablue!28,line width=0.55pt] (-79mm,7mm) -- (79mm,7mm);

  \node[input] (left) at (-42mm,0) {
    \textbf{Left RGB image}\\
    resize to \(224\!\times\!224\); ImageNet normalisation};
  \node[input] (right) at (42mm,0) {
    \textbf{Right RGB image}\\
    resize to \(224\!\times\!224\); ImageNet normalisation};

  \node[encoder] (leftencoder) at (-42mm,-25mm) {
    \textbf{Shared-weight ResNet-18}\\[-0.2ex]
    left view processed as three-channel RGB\\
    ImageNet-pretrained; early stages frozen\\
    \texttt{layer4} fine-tuned $\rightarrow$ 512-D feature \(f_L\)};
  \node[encoder] (rightencoder) at (42mm,-25mm) {
    \textbf{Shared-weight ResNet-18}\\[-0.2ex]
    right view processed as three-channel RGB\\
    ImageNet-pretrained; early stages frozen\\
    \texttt{layer4} fine-tuned $\rightarrow$ 512-D feature \(f_R\)};

  \node[learned] (concat) at (0,-51mm) {
    \textbf{Feature fusion}\qquad
    \([f_L;f_R]\): \(512+512\rightarrow1024\) values};
  \node[learned,minimum height=17mm] (head) at (0,-69mm) {
    \textbf{Post-fusion prediction head}\\[-0.2ex]
    LayerNorm(1024) $\rightarrow$ Linear(128) $\rightarrow$ SiLU
    $\rightarrow$ Dropout(0.05) $\rightarrow$ Linear(200)};

  \node[contract] (contract) at (0,-95mm) {
    \textbf{Fixed named state proposal: 200 normalised outputs}\\[0.6mm]
    Decoder identities: \texttt{L1 L2 L3 L4 R1 R2 R3 R4}\\[0.4mm]
    Per-identity semantics: \textcolor{mfjablue}{18 categorical scores}
    (side 2, block 14, payload 2)\\[-0.1ex]
    + \textcolor{mfjaorange}{7 numeric targets}
    (bounding box 4, \(s_m\), \(s_{ratio}\), auxiliary \(L_{segment}\))\\[0.3mm]
    \textbf{\(8\times(18+7)=200\)} };

  \draw[flow] (left) -- (leftencoder);
  \draw[flow] (right) -- (rightencoder);
  \draw[flow] (leftencoder.south) -- ([xshift=-42mm]concat.north);
  \draw[flow] (rightencoder.south) -- ([xshift=42mm]concat.north);
  \draw[flow] (concat) -- (head);
  \draw[flow] (head) -- (contract);

  % Panel B: the configured registry decides which identity values may be
  % decoded; the neural network itself never decides presence.
  \node[paneltitle,anchor=center] at (0,-117mm)
    {B \enspace RUNTIME DECODING AND CHECKS};
  \draw[mfjablue!28,line width=0.55pt] (-79mm,-121mm) -- (79mm,-121mm);

  \node[runtimeinput] (proposal) at (-40mm,-133mm) {
    \textbf{Model proposal + selected artefacts}\\
    denormalise with training statistics};
  \node[runtimeinput] (presence) at (40mm,-133mm) {
    \textbf{Current identity registry}\\
    \texttt{ShuttleState} names and timestamps only};

  \node[gate] (decode) at (0,-153mm) {
    \textbf{Identity-aware presence gate}\\[-0.2ex]
    pair each fixed decoder identity with its independent
    \texttt{present}, \texttt{absent} or \texttt{unknown} state};

  \node[presentcase] (present) at (-53mm,-177mm) {
    \textbf{present}\\[-0.2ex]
    decode this identity's learned state values};
  \node[absentcase] (absent) at (0,-177mm) {
    \textbf{absent}\\[-0.2ex]
    ignore its learned values; retain configured absence};
  \node[rejectcase] (reject) at (53mm,-177mm) {
    \textbf{unknown / conflict}\\[-0.2ex]
    reject this observation; publish no visual update};

  \node[accepted] (accepted) at (0,-203mm) {
    \textbf{After all identities are consistent: build the visual observation}\\[-0.2ex]
    check artefacts, freshness, vocabulary and numeric consistency\\
    $\rightarrow$ publish \texttt{VisualStateObservation}};

  % Keep every connector vertical by using aligned ports on the wide nodes.
  % This preserves the two independent inputs without diagonal crossings.
  \draw[flow] ([xshift=-40mm]contract.south) -- (proposal.north);
  \draw[gateflow] (proposal.south) -- ([xshift=-40mm]decode.north);
  \draw[gateflow] (presence.south) -- ([xshift=40mm]decode.north);
  \draw[caseflow] ([xshift=-53mm]decode.south) -- (present.north);
  \draw[caseflow] (decode.south) -- (absent.north);
  \draw[caseflow] ([xshift=53mm]decode.south) -- (reject.north);
  \draw[gateflow] (present.south) -- ([xshift=-53mm]accepted.north);
  \draw[gateflow] (absent.south) -- (accepted.north);

  \path[use as bounding box] (-82mm,14mm) rectangle (82mm,-216mm);
\end{tikzpicture}

  \caption[Rail-isolated visual model and deterministic state derivation.]
  {The model shares low-level ResNet-18 weights through \code{layer3}, then keeps
  \code{layer4} and the four-query prediction head independent for each rail.
  The model produces 168 learned values; side, segment length and $s_m$ are
  derived deterministically before whole-observation validation.}
  \label{fig:visual-model-contract}
\end{figure}
\FloatBarrier

Presence is an independent input, not a prediction. Training masks any identity
that is absent or lacks a required target. At runtime, every present identity
must pass artefact, camera-order, freshness, vocabulary, numeric and confidence
checks; any failure rejects the whole observation rather than publishing a
partial state.

\begin{table}[htbp]
\centering
\small
\setlength{\tabcolsep}{4pt}
\caption{Learned objectives and their role in the selected V4 training run.}
\label{tab:visual-training-objectives}
\begin{tabularx}{\textwidth}{@{}
>{\raggedright\arraybackslash}p{0.22\textwidth}
>{\raggedright\arraybackslash}p{0.16\textwidth}
>{\raggedright\arraybackslash}X@{}}
\toprule
\textbf{Output} & \textbf{Loss weight} & \textbf{Role}\\
\midrule
Public segment & Cross-entropy, 4 & Primary discrete location used by symbolic
planning\\ \tablerowrule
Within-segment ratio & Smooth-L1, 2 & Continuous longitudinal refinement within
the selected segment\\ \tablerowrule
Loaded state & Cross-entropy, 1 & Payload constraint used for grounded task
selection\\ \tablerowrule
Bounding box & L1/generalised intersection over union (GIoU)
\citep{giou}, 1 &
Identity-localised visual support and diagnostic overlap\\ \tablerowrule
Topology distance & Graph distance, 0.25 & Penalises mistakes according to their
distance on the known route map\\
\bottomrule
\end{tabularx}
\end{table}

\subsection{Training and checkpoint selection}

The V1--V3 visual-state development cycle was trained on the KAIROS cluster
operated by CALMIP, using an NVIDIA GH200 GPU. The final V4 training,
fine-tuning and evaluation were performed locally on an NVIDIA RTX~3080 Laptop
GPU.

Verified predecessor weights initialised the shared stem through \code{layer3}
and each independent \code{layer4}; no joint head transferred, and rail-branch
parameters, batch-normalisation running statistics (BatchNorm buffers) and
query heads were optimised independently.

The newly initialised query heads were trained before progressively exposing
the deeper rail branches and shared stem. This reduced the chance that an
unstable new head would immediately disturb inherited features. Appearance
augmentation was allowed, whereas camera swaps and side-breaking flips were
forbidden because they would violate the rail contract. The exact optimisation,
sampling and evaluation controls are collected in
\cref{tab:v4-training-controls}.

\begin{table}[H]
\centering
\footnotesize
\renewcommand{\arraystretch}{1.04}
\caption{Selected V4 optimisation, selection and evaluation controls.}
\label{tab:v4-training-controls}
\begin{tabularx}{\textwidth}{@{}
>{\raggedright\arraybackslash}p{0.20\textwidth}
>{\raggedright\arraybackslash}p{0.43\textwidth}
>{\raggedright\arraybackslash}X@{}}
\toprule
\textbf{Control} & \textbf{Frozen setting} & \textbf{Purpose}\\
\midrule
Optimisation & AdamW \citep{adamw}; learning rates $3\times10^{-4}$ for the heads,
$10^{-4}$ for \code{layer4} and $2\times10^{-5}$ for shared \code{layer3};
weight decay $10^{-4}$ & Separate rates protect inherited low-level features\\
\tablerowrule
Staged unfreezing & Heads from epoch~1, independent \code{layer4} from epoch~2
and shared \code{layer3} from epoch~4 & Introduce capacity from task-specific to
shared features\\ \tablerowrule
Epoch definition & 8,000 records sampled with replacement: 4,000 legacy and
4,000 hard-case records; batch 24 and mixed precision & Keep both pools present
while emphasising rare cells and boundaries\\ \tablerowrule
Candidate selection & Planning score below, then worst rail--segment recall,
then lowest 95th-percentile correct-segment position error; patience four & Choose one
checkpoint from Validation using criteria fixed before training\\ \tablerowrule
Evaluation firewall & Seed 31520260811; post-selection regression excluded from
selection and calibration; Final Test unavailable until freeze & Prevent later
evaluation data from influencing the selected model\\
\bottomrule
\end{tabularx}
\end{table}

Macro F1 averages the class-wise harmonic mean of precision and recall, giving
each segment class equal weight. The fixed Validation planning score was
\[
0.30F_{1,\mathrm{macro}}+0.15A_{\mathrm{min\,side}}
+0.20R_{\mathrm{worst\,cell}}+0.20A_{\mathrm{joint},0.05}
+0.075A_{\mathrm{switch}}+0.075A_{\mathrm{boundary}}.
\]
All 12 epochs completed, the procedure selected epoch~11 and the patience
condition was not reached. The regression set could reject this frozen
candidate but could neither choose a checkpoint nor tune it; the Final Test
remained unavailable throughout.

Validation also fitted the segment-score calibration and online acceptance
floor. Loaded-state score remained uncalibrated. The runtime binds checkpoint,
architecture, preprocessing, camera order and topology fingerprint before
becoming ready. \Cref{tab:visual-threshold-map} separates these online checks
from labelled offline tolerances; \cref{app:reproducibility} binds the exact
selected-run artefacts.

\subsection{Evidence after the V3-to-V4 redesign}

The architectural diagnosis in the preceding sections was followed by a large
change on the retained hard-case Validation comparison
(\cref{tab:v3-v4-development-comparison}). The comparison is useful evidence of
the complete redesign and current training protocol, but it is not a controlled
ablation: resizing, branch isolation, output definition, data weighting and
training schedule changed together. It therefore does not prove how much of the
improvement came from any one change.

\begin{table}[htbp]
\centering
\small
\caption{Development comparison before and after the V4 redesign.}
\label{tab:v3-v4-development-comparison}
\begin{tabularx}{\textwidth}{@{}p{0.36\textwidth}p{0.20\textwidth}p{0.20\textwidth}X@{}}
\toprule
\textbf{Hard-case Validation measure} & \textbf{V3} & \textbf{V4} &
\textbf{Change}\\
\midrule
Public-segment top-1 accuracy & 51.6415\% & 99.9113\% & +48.2698 points\\
\tablerowrule
Loaded-state accuracy & 66.1934\% & 99.8669\% & +33.6735 points\\
\bottomrule
\end{tabularx}
\end{table}

The comparison was development evidence: it justified carrying V4, rather than
V3, into the frozen qualification sequence, but it was not counted as Final-Test
performance. I stopped redesigning only after the selected epoch also met the
predeclared Validation and post-selection regression gates and its artefact,
preprocessing and topology identities were frozen. The independent
opposite-camera blank/shuffle check later changed own-side V4 outputs by zero at
recorded precision, supporting the intended rail isolation. Passing these gates
defined the decision to test the model; they were selection controls rather
than a guarantee of future predictions.

\subsection{Preregistered Final Test}

Before evaluation, I froze the checkpoint, scenario plan, evaluator, thresholds
and single-run policy. No Final-Test result could alter training, checkpoint
selection, calibration or thresholds. Here, \emph{preregistered} means that the
repository controls and hashes were fixed before any Test prediction, not that
the protocol was registered externally. The one-shot method, audit and results
are reported in \cref{sec:visual-v4-results}; their records are indexed through
\cref{app:reproducibility}.

The final AI layer emits only a confirmed task and scene record; deterministic
planning and supervision remain responsible for commands. The V3 failures
motivated the rail-isolated V4 design, whose frozen evaluation covers only the
known fleet and fixed Gazebo cameras (\cref{sec:visual-v4-results}).
